% !TEX program = xelatex
% Biên dịch: xelatex 05-tier-a-reading-note.tex
% Tuân thủ 00-methodology.tex

\documentclass[11pt,a4paper]{article}

%=========================================================
% PACKAGES
%=========================================================
\usepackage[margin=2.2cm]{geometry}
\usepackage{fontspec}
\usepackage{polyglossia}
\setdefaultlanguage{vietnamese}
\setotherlanguage{english}
\setmainfont{Latin Modern Roman}
\setsansfont{Latin Modern Sans}
\setmonofont{Latin Modern Mono}

\usepackage{amsmath}
\usepackage{amssymb}
\usepackage{physics}
\usepackage{xcolor}
\usepackage{graphicx}
\usepackage{tabularx}
\usepackage{booktabs}
\usepackage{longtable}
\usepackage{array}
\usepackage{enumitem}
\usepackage{titlesec}
\usepackage{fancyhdr}
\usepackage{hyperref}
\usepackage{bookmark}
\usepackage{url}
\usepackage{tcolorbox}
\tcbuselibrary{skins, breakable}

%=========================================================
% STYLE
%=========================================================
\definecolor{accent}{HTML}{0B5FA5}
\definecolor{accent2}{HTML}{B8860B}
\definecolor{success}{HTML}{2E8B57}
\definecolor{soft}{HTML}{F4F6FA}
\definecolor{softyellow}{HTML}{FFF8E1}
\definecolor{softgreen}{HTML}{E8F5E9}

\hypersetup{colorlinks=true, linkcolor=accent, urlcolor=accent, citecolor=accent2,
  pdftitle={05 - Tier A Reading Notes - QML Foundations}, pdfauthor={TS. Do Phuc Hao}}

\titleformat{\section}{\Large\bfseries\color{accent}}{\thesection.}{0.5em}{}
\titleformat{\subsection}{\large\bfseries\color{accent}}{\thesubsection}{0.5em}{}

\setlist[itemize]{leftmargin=1.2em, itemsep=2pt, topsep=3pt}
\setlist[enumerate]{leftmargin=1.4em, itemsep=2pt, topsep=3pt}

\pagestyle{fancy}
\fancyhf{}
\rhead{\small\textit{Tier A Reading Notes}}
\lhead{\small TS. Đỗ Phúc Hảo}
\rfoot{\thepage}
\renewcommand{\headrulewidth}{0.3pt}

\newtcolorbox{take}[1][]{enhanced, breakable, colback=softgreen, colframe=success,
  boxrule=0.6pt, arc=2pt, left=6pt, right=6pt, top=4pt, bottom=4pt,
  title={#1}, fonttitle=\bfseries\color{white}, coltitle=white, colbacktitle=success}

\newtcolorbox{warn}[1][]{enhanced, breakable, colback=softyellow, colframe=accent2,
  boxrule=0.6pt, arc=2pt, left=6pt, right=6pt, top=4pt, bottom=4pt,
  title={#1}, fonttitle=\bfseries\color{white}, coltitle=white, colbacktitle=accent2}

\newtcolorbox{note}[1][]{enhanced, breakable, colback=soft, colframe=accent,
  boxrule=0.4pt, arc=2pt, left=6pt, right=6pt, top=4pt, bottom=4pt,
  title={#1}, fonttitle=\bfseries\color{white}, coltitle=white, colbacktitle=accent}

%=========================================================
% METADATA
%=========================================================
\title{\vspace{-1.5cm}\textbf{\color{accent} 05 --- Tier A Reading Notes} \\[3pt]
       \Large QML Foundations --- 4 paper phải đọc \\[2pt]
       \normalsize \textit{Biamonte 2017 $\cdot$ Schuld 2021 $\cdot$ Cerezo 2021 $\cdot$ McClean 2018}}
\author{TS. Đỗ Phúc Hảo \\ \small Khoa CNTT --- Đại học Kiến trúc Đà Nẵng}
\date{Ngày 24 tháng 04 năm 2026 (phiên bản 0.1)}

\begin{document}
\maketitle
\thispagestyle{fancy}

\begin{note}[Mục đích]
Tóm lược 4 foundation paper (Tier A ở file 01) với \textbf{đủ sâu} để tác giả có thể tự viết section ``Preliminaries'' trong paper của mình. Mỗi note gồm: (1) bối cảnh \& động cơ, (2) ý tưởng cốt lõi + công thức chính, (3) kết quả nổi bật, (4) \textit{takeaway cho dự án DynaQGNN-IDS}, (5) câu hỏi mở / hướng research trực tiếp nối dài.
\end{note}

\tableofcontents
\vspace{0.4cm}
\hrule
\vspace{0.4cm}

%=========================================================
\section{Biamonte et al. (2017) --- ``Quantum Machine Learning''}
%=========================================================

\textbf{Venue}: \textit{Nature}, vol. 549, pp. 195--202, 2017. Ref: Biamonte, Wittek, Pancotti, Rebentrost, Wiebe, Lloyd.

\subsection{Bối cảnh \& động cơ}

Paper survey đầu tiên trên tạp chí top định hình lĩnh vực QML. Ra đời khi D-Wave quantum annealer đã hoạt động, NISQ era chưa bắt đầu. Câu hỏi trung tâm: \textit{quantum computing có thể tăng tốc ML như thế nào?}

\subsection{Ý tưởng cốt lõi}

Ba nguồn ``quantum advantage'' khả dĩ cho ML:

\begin{enumerate}
  \item \textbf{Speedup thuật toán tuyến tính}:
    Thuật toán HHL (Harrow, Hassidim, Lloyd, 2009) giải hệ $\mathbf{A}\mathbf{x} = \mathbf{b}$ trong $\mathcal{O}(\log N \cdot \kappa^2 / \epsilon)$ thay vì $\mathcal{O}(N)$ cổ điển (có điều kiện: $\mathbf{A}$ sparse + well-conditioned + chỉ đọc few amplitudes của $\mathbf{x}$).
  \item \textbf{Mở rộng không gian đặc trưng qua feature map lượng tử}:
    Encoding $\mathbf{x} \mapsto \ket{\phi(\mathbf{x})}$ cho phép biểu diễn trong không gian Hilbert $2^n$-chiều. Kernel $K(\mathbf{x}, \mathbf{x}') = |\braket{\phi(\mathbf{x})|\phi(\mathbf{x}')}|^2$.
  \item \textbf{Sampling từ phân phối khó}:
    Quantum Boltzmann machines sample từ Gibbs distribution trên model lượng tử.
\end{enumerate}

\paragraph{Taxonomy QML 2$\times$2}:

\begin{longtable}{@{}p{2.5cm} p{5.5cm} p{6cm}@{}}
\toprule
 & \textbf{Data: Classical} & \textbf{Data: Quantum} \\
\midrule
\endhead
\textbf{Algo: Classical} & Standard ML (baseline) & Analyze quantum experiments by classical ML \\
\textbf{Algo: Quantum} & QML cho classical data (MAJORITY of papers) & Quantum ML cho quantum data (e.g., Hamiltonian learning) \\
\bottomrule
\end{longtable}

\subsection{Kết quả nổi bật}
\begin{itemize}
  \item Chứng minh HHL-based QSVM, Quantum PCA, Quantum clustering với exponential speedup lý thuyết.
  \item Cảnh báo về \textbf{caveat data loading}: nếu phải encode classical data vào qubit từng bit thì speedup bị triệt tiêu $\Rightarrow$ cần QRAM (không khả thi hardware).
  \item Chỉ ra quantum annealing có thể dùng cho sampling model-based.
\end{itemize}

\begin{take}[Takeaway cho DynaQGNN-IDS]
\begin{itemize}
  \item \textbf{Encoding đúng là then chốt}: Paper của ta phải so sánh \textit{angle} vs \textit{amplitude} vs \textit{re-uploading} và giải thích lý do chọn. Đây là gốc rễ của quantum advantage thật hay ảo.
  \item \textbf{Không hứa hẹn exponential speedup}: với NISQ, ta nhắm ``representation efficiency'' (ít qubit hơn để đạt expressivity tương đương classical depth) và ``noise resilience'' thay vì raw speedup.
  \item \textbf{Phải trả lời caveat data loading}: paper của ta load đồ thị NTN snapshot (adjacency + node features) --- phải thảo luận cost encoding trong Complexity Analysis section.
\end{itemize}
\end{take}

\subsection{Câu hỏi mở}
\begin{itemize}
  \item Khi nào kernel trick lượng tử có ``provable advantage''? (đã được Havlíček 2019 + Huang 2021 trả lời một phần).
  \item NISQ-friendly QML algorithms? $\rightarrow$ VQA (Cerezo 2021 xuống sau).
\end{itemize}

%=========================================================
\section{Schuld \& Petruccione (2021) --- ``ML with Quantum Computers''}
%=========================================================

\textbf{Venue}: Springer textbook, 2nd edition (mở rộng từ 2018 edition).

\subsection{Bối cảnh \& động cơ}

Sách giáo khoa chuẩn đầu tiên cho QML dành cho cả physicist và ML practitioner. Biên soạn kịp thời với sự bùng nổ NISQ / PennyLane (Schuld là co-founder của Xanadu PennyLane).

\subsection{Cấu trúc \& ý tưởng cốt lõi}

\paragraph{Phần I --- Foundations}: Cơ sở QC (Hilbert, unitary, measurement) + ML (supervised, unsupervised). Dirac notation bị lấy làm language chính.

\paragraph{Phần II --- Information Encoding}: Chi tiết \textbf{4 loại encoding}:
\begin{enumerate}
  \item \textbf{Basis encoding}: mỗi bit classical $\mapsto$ 1 qubit (cost: $N$ qubit cho $N$-bit data).
  \item \textbf{Amplitude encoding}: vector $\mathbf{x} \in \mathbb{R}^{2^n}$ $\mapsto \ket{\psi} = \sum_i x_i \ket{i}/\|\mathbf{x}\|$; chỉ cần $n = \log_2 N$ qubit. \textbf{Nhưng} cost chuẩn bị state $\mathcal{O}(N)$ gate.
  \item \textbf{Angle encoding}: $\mathbf{x} \mapsto \bigotimes_i R_y(\theta_i) \ket{0}$; đơn giản, hiệu quả cho NISQ.
  \item \textbf{Hamiltonian / Data re-uploading}: encode data nhiều lần vào circuit $\rightarrow$ tăng expressivity (Pérez-Salinas 2020).
\end{enumerate}

\paragraph{Phần III --- Quantum Models}:
\begin{itemize}
  \item Variational quantum circuits (VQC) như \textit{parametrized function class}.
  \item Kernel methods: quantum kernel estimation, chuyển QSVM thành kernel-based classifier.
  \item Born machines: generative model dựa trên Born rule.
\end{itemize}

\paragraph{Phần IV --- Training}:
\begin{itemize}
  \item \textbf{Parameter-shift rule}: gradient của $\braket{\hat{H}}_{\boldsymbol{\theta}}$ tính được bằng cách chạy circuit với $\theta_k \pm \pi/2$ (không cần adjoint differentiation phức tạp).
  \item \textbf{Analytic vs hardware gradients}: trên QPU thật, cần tối thiểu 2 circuit evaluations per parameter per step.
  \item \textbf{Noise + shot noise}: gradient estimator có variance $\propto 1/S$ (shots); thực tế cần thousand shots.
\end{itemize}

\begin{take}[Takeaway cho DynaQGNN-IDS]
\begin{itemize}
  \item \textbf{Encoding strategy cho đồ thị NTN}: adjacency $\mathbf{A}_t \in \{0,1\}^{K \times K}$ + node features $\mathbf{X}_t$. Đề xuất: \textit{hybrid encoding} --- angle cho node features, amplitude/block-encoding cho adjacency (theo Lloyd 2016).
  \item \textbf{Parameter-shift rule + PennyLane}: triển khai chuẩn, không tự viết. Tận dụng \texttt{pennylane-lightning-gpu} với adjoint differentiation cho simulator.
  \item \textbf{Shot noise trong real QPU}: khi chạy IBM hardware validation cho paper B, cần budget $\geq 8000$ shots/circuit.
\end{itemize}
\end{take}

\subsection{Chương cần đọc kỹ}
\begin{enumerate}
  \item Chương 5 (Encoding) --- nền cho thiết kế DynaQGNN encoder.
  \item Chương 8 (Training) --- parameter-shift, barren plateau intro.
  \item Chương 9 (Quantum kernel methods) --- vì QSVM là baseline.
\end{enumerate}

%=========================================================
\section{Cerezo et al. (2021) --- ``Variational Quantum Algorithms''}
%=========================================================

\textbf{Venue}: \textit{Nature Reviews Physics}, vol. 3, pp. 625--644, 2021. Ref: Cerezo, Arrasmith, Babbush, Benjamin, Endo, Fujii, McClean, Mitarai, Yuan, Cincio, Coles.

\subsection{Bối cảnh \& động cơ}

Survey định hình ``VQA era''. Trả lời câu hỏi: trong NISQ thời đại, thuật toán lượng tử nào có cơ may thắng? Trả lời: \textbf{variational hybrid}.

\subsection{Khung VQA tổng quát}

Mỗi VQA được đặc trưng bởi 4 thành phần:

\begin{enumerate}
  \item \textbf{Cost function} $C(\boldsymbol{\theta})$: đặc trưng bài toán, typically $C(\boldsymbol{\theta}) = \braket{\psi(\boldsymbol{\theta})|\hat{H}|\psi(\boldsymbol{\theta})}$.
  \item \textbf{Ansatz} $U(\boldsymbol{\theta})$: parameterized circuit, chia 2 nhóm:
    \begin{itemize}
      \item Problem-inspired (VQE với chemistry ansatz, QAOA layers).
      \item Hardware-efficient (HEA) --- đơn giản, native gates.
    \end{itemize}
  \item \textbf{Gradient computation}: parameter-shift, adjoint, SPSA (shot-efficient).
  \item \textbf{Classical optimizer}: Adam, L-BFGS, SPSA, Nelder-Mead.
\end{enumerate}

\paragraph{Các flagship VQA}:
\begin{itemize}
  \item \textbf{VQE} (Peruzzo 2014) --- tìm ground state energy, chemistry.
  \item \textbf{QAOA} (Farhi 2014) --- combinatorial optimization (MaxCut, portfolio\dots).
  \item \textbf{VQC / QNN}: classification, regression.
  \item \textbf{VQLS} (Variational Quantum Linear Solver) --- alternative cho HHL NISQ-friendly.
  \item \textbf{QGAN}: generative.
\end{itemize}

\subsection{Thách thức chính}

\begin{warn}[Ba thách thức khi áp VQA]
\begin{enumerate}
  \item \textbf{Barren plateau} (McClean 2018): gradient $\mathrm{Var}[\partial_\theta C] \propto 2^{-n}$ cho HEA sâu $\Rightarrow$ training thất bại.
  \item \textbf{Expressibility vs trainability trade-off}: ansatz ``expressive'' thường khớp với barren plateau; ansatz ``trainable'' thường kém biểu đạt.
  \item \textbf{Noise}: quantum noise làm landscape phẳng hơn, tăng barren plateau (Wang 2021).
\end{enumerate}
\end{warn}

\subsection{Chiến lược đối phó (reviewer BẮT BUỘC muốn thấy)}
\begin{itemize}
  \item \textbf{Shallow circuits + layer-wise training} (Grant 2019; Skolik 2021).
  \item \textbf{Local cost functions} (Cerezo 2021a): đo từng qubit thay vì global observable.
  \item \textbf{Smart initialization}: Gaussian / identity block initialization.
  \item \textbf{Problem-inspired ansatz}: match symmetry bài toán, loại bỏ parameter dư.
  \item \textbf{Entanglement-aware design}: tránh quá entangled từ đầu.
\end{itemize}

\begin{take}[Takeaway cho DynaQGNN-IDS]
\begin{itemize}
  \item \textbf{Chọn ansatz}: đề xuất \textit{topology-aware problem-inspired} --- ansatz mirror cấu trúc adjacency $\mathbf{A}_t$ (CNOT pattern = edges); không dùng HEA thuần. Điều này cũng giúp né barren plateau.
  \item \textbf{Parameter-shift} + \textbf{adjoint differentiation} trên simulator; khi validation trên IBM hardware thì chỉ parameter-shift.
  \item \textbf{Local cost}: thiết kế detection loss dạng local (per-flow decision) thay vì global (batch-level) để mitigate barren plateau.
  \item \textbf{Phải có section ``Complexity and Trainability Analysis''} trong paper --- reviewer sẽ đòi.
\end{itemize}
\end{take}

%=========================================================
\section{McClean et al. (2018) --- ``Barren Plateaus''}
%=========================================================

\textbf{Venue}: \textit{Nature Communications}, vol. 9, 4812, 2018. Ref: McClean, Boixo, Smelyanskiy, Babbush, Neven.

\subsection{Bối cảnh}

Paper ``báo tử''. Cho thấy rằng naively tăng circuit depth sẽ gặp vấn đề cực nghiêm trọng. Được cite $>$ 2000 lần.

\subsection{Kết quả chính}

Cho ansatz $U(\boldsymbol{\theta}) = \prod_l U_l(\theta_l) V_l$ với $U_l$ random unitary tạo từ gate chạy trên 2-design:

\[
\mathbb{E}[\partial_\theta C] = 0, \quad \mathrm{Var}[\partial_\theta C] = \Theta(2^{-n}).
\]

Nghĩa là với $n$ qubit và ansatz đủ sâu, gradient \textbf{giảm mũ}; với $n = 30$ thì gradient $\sim 10^{-9}$, optimizer nào cũng dừng.

\subsection{Trực giác}

Với unitary đủ random (``tiệm cận Haar measure''), landscape tiệm cận phẳng gần 0. Training khởi tạo random $\Rightarrow$ gradient vanish $\Rightarrow$ không học được.

\subsection{Những gì làm tăng barren plateau}
\begin{enumerate}
  \item Circuit quá sâu.
  \item Ansatz quá expressive.
  \item Global cost function (đo trên nhiều qubit cùng lúc).
  \item Noise tăng (Wang 2021 bổ sung).
  \item Entanglement quá cao từ đầu.
\end{enumerate}

\subsection{Những gì giảm barren plateau}
\begin{enumerate}
  \item Shallow circuit ($\leq \mathcal{O}(\log n)$ depth).
  \item Problem-inspired ansatz, structured.
  \item Local cost function.
  \item Smart initialization (identity block, Gaussian small).
  \item Layer-wise / block-wise training.
  \item Pre-training với smaller task.
\end{enumerate}

\begin{warn}[Quy tắc kiểm tra bắt buộc cho paper]
Mọi paper VQC/QNN phải báo cáo $\mathrm{Var}[\partial_\theta C]$ theo số qubit $n \in \{4, 8, 12, 16\}$. Nếu thấy giảm mũ $\rightarrow$ thừa nhận \& mitigate. Nếu không $\rightarrow$ giải thích tại sao.
\end{warn}

\begin{take}[Takeaway cho DynaQGNN-IDS]
\begin{itemize}
  \item Bắt buộc có \textbf{barren plateau diagnostic}: plot $\mathrm{Var}[\partial_\theta C]$ vs $n$. Đây là phần của ablation A5 (xem file 04).
  \item Phải có \textbf{ablation ``BP mitigation enabled vs disabled''} để chứng minh tác dụng của kỹ thuật mình chọn.
  \item \textbf{Chọn depth} không quá sâu: mục tiêu $D \in \{2, 4, 6\}$ layers; $D = 8$ chỉ để cho thấy ``mới bắt đầu bị BP''.
  \item \textbf{Local cost}: thiết kế loss per-node hoặc per-edge, không tổng qubit cùng lúc.
\end{itemize}
\end{take}

%=========================================================
\section{Ghép nối 4 paper thành 1 section ``Preliminaries'' cho DynaQGNN-IDS}
%=========================================================

Dưới đây là skeleton mà tác giả có thể dùng trực tiếp cho Section II của paper A:

\begin{note}[Skeleton Section II --- Preliminaries]
\textbf{II.A Quantum Machine Learning Primer}: dùng Biamonte 2017 để giới thiệu 3 nguồn quantum advantage, nhấn mạnh feature map và kernel methods.

\textbf{II.B Variational Quantum Circuits}: dùng Schuld 2021 + Cerezo 2021 để định nghĩa VQC, ansatz (HEA vs problem-inspired), parameter-shift rule, hybrid optimization loop.

\textbf{II.C Encoding Classical Graph Data}: Schuld 2021 Chương 5; cụ thể cho ta: angle encoding cho node features + block-encoding adjacency (mention Lloyd 2016 graph state preparation).

\textbf{II.D Barren Plateaus and Mitigation}: McClean 2018; trình bày ngắn công thức $\mathrm{Var}[\partial_\theta C] = \Theta(2^{-n})$; liệt kê 5 chiến lược mitigation; khẳng định DynaQGNN dùng 3 trong số đó (shallow + problem-inspired + local cost).
\end{note}

%=========================================================
\section{Danh sách đọc nối dài ngay sau Tier A}
%=========================================================

Sau khi nắm Tier A, đọc ngay 5 paper ngắn sau để hoàn thiện nền:

\begin{enumerate}
  \item \textbf{Schuld \& Killoran 2019} --- ``Quantum ML in Feature Hilbert Spaces'' --- kernel view chặt chẽ.
  \item \textbf{Havlíček et al. 2019} --- ``Supervised learning with quantum-enhanced feature spaces'' (Nature) --- quantum kernel thực sự thắng khi nào.
  \item \textbf{Mitarai et al. 2018} --- ``Quantum circuit learning'' --- định nghĩa parameter-shift rule.
  \item \textbf{Pérez-Salinas et al. 2020} --- ``Data re-uploading for universal approximation'' --- trick tăng expressivity ansatz nông.
  \item \textbf{Huang et al. 2021} --- ``Power of data in QML'' (Nature Communications) --- chỉ ra khi nào classical + data thắng quantum.
\end{enumerate}

%=========================================================
\section{Tiếp theo}
%=========================================================

Sau file này, theo kế hoạch sẽ sang:
\begin{itemize}
  \item \textbf{figures/fig-system-model.tex} --- TikZ standalone system model SAGIN.
  \item \textbf{figures/fig-dynaqgnn-ansatz.tex} --- TikZ circuit diagram DynaQGNN.
  \item \textbf{Repo skeleton} \texttt{dynaqgnn-ntn-ids/}.
\end{itemize}

\vspace{0.4cm}
\hrule
\vspace{0.3cm}
\textit{\small Phiên bản 0.1 --- 24/04/2026. Sau khi đọc đủ 4 paper + 5 paper nối dài, bump lên 0.2 với gạch chân ý quan trọng được highlight khi đọc.}

\end{document}
